\subsection{QEMU Device Model (Qdev)}
\label{chap:qemu_qdev}

While the QOM provides the infrastructure for the creation and the handling 
of objects in QEMU, these are generic and they do not specifically address the 
requirements of hardware device modeling. As result, QEMU Device API (QDev) 
was designed to covered this gap. This QEMU's API is a specialized framework built 
atop QOM because of this, each QDev device is itself a QOM object, meaning it inherits 
QOM's type system, introspection, and dynamic property mechanisms.

Qdev functions as an abstraction layer for the device lifecycle, unifying the 
description, creation, connection, and configuration of the emulated hardware components,
and ensuring that every device, regardless of its specific peripheral function, follows consistent patterns for initialization, 
property management, lifecycle stages, and system integration. The device-specific 
capabilities that Qdev introduce for hardware emulation includes: DeviceClass, 
Lifecycle Stages, Bus Abstraction, GPIO/IRQ System, Reset Propagation, Hotplug Support.

Most of the capabilities mentioned are achieved through several layers of 
abstraction in order to standardize and simplify the creation of emulated hardware to 
the developer. Due to the associated complexity, it is important for developers to 
understand qdev's layers of abstraction before reading or writing code; not doing it, 
could lead to a spiral of confusion about how and why things are done in QEMU. 
The device lifecycle is the first thing Qdev applies, providing a guide with several 
standardized phases to ensure that the firmware running on QEMU can reliably interact 
with devices through predictable boot sequences and initialization patterns. 
The **F1** show this process.

- **F1**

The second element provided by QDev is a bus abstraction layer, which will be 
especially relevant later in this thesis for emulating the STM32F405 address map. 
In QEMU, thanks to this bus abstraction layer, devices do not exist in isolation 
but are organized hierarchically and connected to buses. These buses manage 
address allocation, establish parent-child relationships, and handle address translation 
for memory-mapped I/O. This is fundamental for emulating microcontrollers with 
multiple peripherals, each occupying different memory regions and responding to 
specific addresses on the bus. Table T1 shows a list of common QEMU buses 
and their purpose.

- **T1**

Fourth, QDev implements reset propagation, allowing a system reset command 
to cascade through the entire device tree, ensuring all peripherals return to their 
initial states in the correct order—essential for firmware that relies on peripheral 
reset behavior. And finally, QDev also supports hotplug, permitting devices to be 
dynamically added or removed at runtime, though this feature is less critical for 
bare-metal microcontroller emulation than for system-level scenarios.

Like the QOM API, QDev also has its main types of C structures. There are 
three types of C structures that developers should be familiar with: **DeviceClass**, 
**DeviceState**, and **BusState**. Of these, the one developers should be familiar 
with is **DeviceClass**, since all device types inherit from **DeviceClass**, which 
in turn extends **ObjectClass**.

-**C2**